\cvsection{Expériences professionnelles}

\begin{cventries}
  \cventry
    {Stagiaire ingénieur de recherche en IA}
    {INRIA}
    {Bordeaux, France}
    {Mars 2026 -- Septembre 2026}
    {
      \begin{cvitems}
        \item {Analyse du besoin d’aider les chercheurs à identifier des articles scientifiques pertinents et définition d’une solution adaptée à leurs usages.}
        \item {Conception en autonomie d’une application intégrant un agent IA multi-étapes, depuis le cadrage fonctionnel jusqu’à la réalisation d’un prototype utilisable.}
        \item {Traduction du cas d’usage en composants techniques : agent, prompts, outils, gestion du contexte, architecture RAG, API et interface web.}
        \item {Développement d’un backend Python/FastAPI, d’API REST documentées avec OpenAPI, d’une persistance MongoDB et d’une interface React/Next.js.}
        \item {Mise en œuvre d’un pipeline RAG avec chunking, embeddings, recherche sémantique et base vectorielle ChromaDB.}
        \item {Construction d’un protocole d’évaluation et comparaison de plusieurs LLM et quantifications selon la qualité, la latence, le TTFT et la taille du contexte.}
        \item {Analyse et synthèse des résultats afin d’identifier les compromis entre performance, ressources mobilisées et besoins du cas d’usage.}
        \item {Documentation des travaux et rédaction en cours d’un article scientifique présentant la démarche, les résultats et leurs limites.}
      \end{cvitems}
    }

  \cventry
    {Stagiaire ingénieur de recherche en apprentissage par renforcement}
    {Université Karunya}
    {Coimbatore, Inde}
    {Juin 2025 -- Sept. 2025}
    {
      \begin{cvitems}
        \item {Développement d’un modèle d’apprentissage par renforcement pour la navigation autonome d’un robot hospitalier, puis déploiement et tests sur Raspberry Pi dans un environnement international.}
      \end{cvitems}
    }
\end{cventries}
